% ============================================================================
% Section 1 -- introduction.
%
% REFRAMED 2026-09-06 (Ben): the paper is a HETEROGENEOUS MEMORY SYSTEM for
% decode, not "a policy layer on top of somebody's kernel". The old framing was
% structurally subordinate -- it invited "so you wrote the scheduler DirectKV
% skipped" -- and it scattered our pieces as a bag of policies. Under this frame
% the allocator, placement, admission, migration and reclaim are components of
% one system that manages HBM and host DRAM as a single KV memory, and the
% capacity model is that system's design rule.
%
% ALSO CUT 2026-09-06: the losslessness material. Exact placement is what moving
% bytes does; arguing for it invites the doubt it was meant to settle, and no
% comparable systems paper does it.
% ============================================================================

\section{Introduction}
\label{sec:intro}

A GPU serving LLM decode requests is mostly idle. Generating one token for a
batch requires reading every model weight and every resident KV cache byte,
with almost no arithmetic per byte read, so the arithmetic units run at
single-digit utilization while the memory system saturates. What prevents the
server from batching more requests --- which is what would put those units to
work --- is not compute. It is that the KV cache has exhausted device memory. A
recent trace-driven study of production deployments reaches the same conclusion
from the workload side: what binds is ``VRAM capacity, not token
budget''~\cite{bottlenecks}.

The machine, however, has more memory than the GPU's. A Grace-Hopper node pairs
96\,GB of HBM with hundreds of gigabytes of host DRAM, and NVLink-C2C connects
them at 450\,GB/s per direction.\footnote{The commonly quoted 900\,GB/s is
bidirectional; a KV stream is unidirectional host-to-device, so 450\,GB/s is
the figure that bounds it, and we measure rather than assume even that.} At
that rate the KV a decoding request needs can cross the link inside the window
during which the GPU is reading its own weights and wants nothing from the link
at all. Host DRAM stops being a place to park idle requests and becomes usable
memory for running ones.

\paragraph{Two memories, managed as one.} This paper builds a
\textbf{heterogeneous memory system for LLM decode}: HBM and host DRAM managed
as a single KV memory. Its load-bearing decision is \emph{how many} requests to
serve from the host tier, and the answer comes from the interconnect rather
than from memory. A host-tier request must move its entire KV across the link
once per decoded token, so the number that can be sustained is set by what the
link delivers inside a decode step --- a quantity this system measures rather
than declares. Around that decision sit the mechanisms it requires: an
allocator in which a host-tier request holds no device blocks, per-request
placement, and migration between tiers under load.

That framing is not how the field currently treats the second memory. Existing
systems use one memory or the other, statically and totally. Stock vLLM and
SGLang keep 100\% of KV in HBM; when it does not fit, requests wait. DirectKV,
which contributes the zero-copy attention kernel that makes host-resident KV
efficient at all, keeps 100\% on the host for every request always --- its \S6
stores each layer's initial $K,V$ in CPU memory during prefill and its Alg.~3
appends new $K,V$ there during decode. That is precisely why it needs no
memory management: placement is static and total, and its \S8.1 says the design
``does not change scheduling, batching, or KV-management policies.''

Neither is optimal, for a reason visible without measurement. An all-HBM server
reads KV at HBM bandwidth and leaves the interconnect idle. An all-host server
reads KV over the interconnect and leaves the HBM not holding weights idle. A
system that places some requests in each reads from both at once, and its batch
is bounded by the sum:

\begin{table}[t]
\centering
\small
\begin{tabular}{lrr}
\toprule
\emph{Projected.} Qwen2.5-32B fp16, 8K, GH200-96 & KV read per step & batch \\
\midrule
all-HBM (vLLM, SGLang)      & 26.5\,GB HBM, \textbf{0 link} & 12 \\
all-host (uniform offload)  & \textbf{0 HBM}, 8.8\,GB link & 4 \\
\textbf{both (this paper)}  & 26.5\,GB HBM \textbf{+} 8.8\,GB link & \textbf{16} \\
\bottomrule
\end{tabular}
\caption{Each single-memory design leaves one of the two KV paths unused.
Weight traffic is identical in all three rows and excluded, so the column
compares like with like. \textbf{These are projections}, computed by
\texttt{model/killer\_v2.py} from: measured checkpoint $W = 65.5$\,GB, a 4\,GB
device reserve for activations and fragmentation, hence $\Ckv = 26.5$\,GB;
$\btok = 256$\,KiB so a request at 8K holds 2.15\,GB; $\mathrm{ITL}_0 = 33.9$\,ms
from Eq.~\ref{eq:step}; $\Xlink = 419$\,GB/s, the measured host-to-device rate
on this part~\cite{pie}, not the 450\,GB/s spec figure; and $\effo = 0.62$.
The last is the weakest input: it is measured on A100/PCIe
(\S\ref{sec:eval:cost}) and assumed to carry to a coherent link, which is the
first thing \S\ref{sec:eval} asks of one. At $\effo = 1$ the host row would
read 6 and the total 18.}
\label{tab:endpoints}
\end{table}

Writing that sum out states the contribution exactly. For requests of $S$
tokens at $\btok{}$ KV bytes each, the resident cohort is a
\emph{memory} quantity and the host cohort is a \emph{bandwidth} one:
%
\begin{equation}
B \;=\; \underbrace{\frac{\Ckv}{S\,\btok{}}}_{\text{device capacity}}
\;+\;
\underbrace{\frac{\mathrm{ITL}\cdot\Xlink\cdot\effo}{S\,\btok{}}}_{\text{interconnect}}
\label{eq:two-terms}
\end{equation}
%
The second term does not contain $\Ckv$. Admission is thereby decoupled from
device capacity and bounded instead by what the link can deliver within a step
--- which is the decision this paper is about, and the one no existing system
makes: the endpoints in Table~\ref{tab:endpoints} set the host cohort to zero
or to everything, neither of which consults the interconnect at all.

\paragraph{When is a second memory tier worth having?} Before building one we
ask when it pays, and the answer is sharply bounded. We derive a closed-form
capacity multiplier $1 + (\Xlink/\BW)\,\effo\,(\HBM/\Ckv)$, where $\Ckv$ is the
device memory left after weights and $\effo$ is the fraction of a rider's
transfer that overlaps compute, and validate its structure against 54 measured
configurations. Two consequences contradict the obvious intuition. Long context
does \emph{not} help: it shrinks the HBM-resident cohort and the affordable
host cohort in the same proportion, leaving their ratio unchanged. And the
value of a second tier tracks \textbf{how much of memory the weights ate},
not model size. An 8B model leaves so much HBM free that there is nothing to
buy; a 32B model at fp16 leaves 28\,GB and gains materially. The sharpest
demonstration is that \emph{the same model at two weight precisions} lands at
two different points on the curve.

The model earns its place by explaining a published result its own authors do
not. DirectKV reports being ``slightly slower than SGLang'' on Llama-3.1-8B and
winning on OPT-13B/30B ``where SGLang OOMs''. Our capacity multiplier predicts
that inversion from the two models' weight footprints alone. A design rule that
accounts for when the technique \emph{fails} is worth more than one tuned to
its successes.

\paragraph{Contributions.}
\begin{enumerate}
\item \textbf{A capacity model for heterogeneous KV memory}
      (\S\ref{sec:model}): a closed-form multiplier and an applicability gate,
      validated on hardware and against a published system's own results. Both
      exclude configurations the field currently evaluates on.
\item \textbf{A two-tier KV memory manager} (\S\ref{sec:design}) in which a
      host-tier request holds no device blocks at all, so tiered requests are
      additive rather than substitutional --- the property that makes the
      interior point of Table~\ref{tab:endpoints} reachable.
\item \textbf{Admission decoupled from device capacity and bounded by the
      interconnect} (\S\ref{sec:design}): the host cohort of
      Eq.~\ref{eq:two-terms} sized from a per-step transfer budget that is
      \emph{measured} --- from the rider-free step time and the achieved link
      bandwidth --- rather than declared as a constant. We show by ablation
      that this derived count, and not a fixed offload fraction, is what
      carries the result.
\item \textbf{An implementation in vLLM v1 and an evaluation separating the
      memory system's contribution from the kernel's} (\S\ref{sec:eval}), by
      running uniform offload and our manager on one substrate, together with
      two negative results we report rather than bury: per-request placement
      driven by observable state did not separate from arrival order on our
      substrate --- and context length, the intuitive signal, does not
      discriminate at all --- while closed-loop migration measured its own
      demotions as harmful and reverted them.
\end{enumerate}

We state the last of these deliberately. A host tier admits a policy design
space --- which requests to move, and when to move them --- and we built that
space and report that, at the rider counts a PCIe-class link permits, it does
not yet pay. A policy choosing among two candidates cannot demonstrate value;
whether it does so when the budget admits dozens is a question for a coherent
interconnect, not a claim we make here.

\paragraph{Scope.} We target a decode node in a prefill/decode-disaggregated
deployment~\cite{distserve,splitwise,mooncake}, where a request arrives with
its KV already computed. \S\ref{sec:background:pd} explains why that assumption
is what makes host-tier requests additive.

\paragraph{Results.} On A100/PCIe --- the regime our own model says cannot pay
--- the memory system works and the economics do not. Against a fixed KV pool,
host-tier requests raise peak concurrency from 8 to 48, itself a lower bound
since 48 was the entire offered load.
% src: paper/data/preload_capacity.json @ f3df05bdb
In a separate experiment that isolates a rider's cost --- riders added on top of
a fixed resident batch, with prefill cancelled by construction --- only
$\effo = 0.62$ of a host-tier request's transfer overlaps compute, and the
unhidden residual costs $8.3$\,ms of a $54$\,ms step, leaving decode throughput
$0.86\times$ that of not offloading at all.
% src: paper/data/rider_cost_4p2b.json @ 21ea9b18d
End to end, through the prefill handover, the deficit is $2.3\%$.
% src: paper/data/pd_decode_goodput.json @ f3df05bdb
That is what the model predicts for a 25\,GB/s link, and DirectKV reaches the
same boundary from the kernel side, concluding that on PCIe theirs ``should be
viewed primarily as a capacity-extension mechanism''. It fixes one end of the
bandwidth axis.
On a coherent interconnect we report the same comparison at matched latency
targets, together with the second point on the bandwidth axis that turns the
cost decomposition of \S\ref{sec:eval:cost} into a law with a measured
crossing rather than an extrapolated one.
